% ===== PRÉAMBULE CANONIQUE — à coller VERBATIM en tête de CHAQUE .tex =====
\documentclass[11pt,a4paper]{article}
\usepackage[utf8]{inputenc}\usepackage[T1]{fontenc}\usepackage{lmodern}
\usepackage[french]{babel}
\usepackage{amsmath,amssymb,mathtools}\usepackage{icomma}
\usepackage{siunitx}\sisetup{locale=FR}  % \qty{}{}, \si{}, \ang{}
\usepackage{xcolor}\usepackage{colortbl}
\usepackage{tikz}\usepackage{pgfplots}\pgfplotsset{compat=1.18}
\usepackage[siunitx,europeanresistors]{circuitikz} % circuits (élec.)
\usepackage[most]{tcolorbox}\usepackage{enumitem}\usepackage{array,booktabs}
\usepackage[a4paper,margin=2.2cm]{geometry}
\usetikzlibrary{arrows.meta,calc,angles,quotes,positioning,patterns,%
  backgrounds,shapes.geometric,decorations.pathmorphing,decorations.markings,intersections}
\definecolor{primary}{RGB}{222,49,99}\definecolor{primarydark}{RGB}{150,22,55}
\definecolor{defcol}{RGB}{0,114,178}\definecolor{methcol}{RGB}{52,92,148}
\definecolor{excol}{RGB}{0,135,90}\definecolor{attncol}{RGB}{200,40,55}
\tcbset{boxbase/.style={enhanced,breakable,boxrule=0.9pt,arc=2.5pt,
  left=3mm,right=3mm,top=2.3mm,bottom=2.3mm,fonttitle=\bfseries,coltitle=white}}
% --- boîtes TUTORIEL (noms EXACTS, ne pas renommer) ---
\newtcolorbox{cle}[1][]{boxbase,colback=defcol!6,colframe=defcol,
  title={\textsc{À retenir}\ifx\relax#1\relax\relax\else\textmd{ : \itshape #1}\fi}}
\newtcolorbox{methode}[1][]{boxbase,colback=methcol!7,colframe=methcol,sharp corners,
  title={\textsc{Méthode}\ifx\relax#1\relax\relax\else\textmd{ --- \itshape #1}\fi}}
\newtcolorbox{exemplebox}[1][]{boxbase,colback=excol!5,colframe=excol,
  title={\textsc{Exemple}\ifx\relax#1\relax\relax\else\textmd{ : \itshape #1}\fi}}
\newtcolorbox{piege}[1][]{boxbase,colback=attncol!6,colframe=attncol,
  title={\textsc{Piège}\ifx\relax#1\relax\relax\else\textmd{ : \itshape #1}\fi}}
\newtcolorbox{remarque}[1][]{boxbase,colback=black!4,colframe=black!45,
  title={\textsc{Remarque}\ifx\relax#1\relax\relax\else\textmd{ : \itshape #1}\fi}}
% --- boîtes EXERCICE / CORRIGÉ (noms EXACTS, auto-comptées) ---
\newtcolorbox[auto counter]{exo}[1][]{boxbase,colback=primary!4,colframe=primary,
  title={\textsc{Exercice}~\thetcbcounter\ifx\relax#1\relax\relax\else\textmd{ --- \itshape #1}\fi}}
\newtcolorbox[auto counter]{corrige}[1][]{boxbase,colback=excol!4,colframe=excol,
  title={\textsc{Corrigé}~\thetcbcounter\ifx\relax#1\relax\relax\else\textmd{ --- \itshape #1}\fi}}
\newcommand{\vv}[1]{\vec{#1}}\newcommand{\ds}{\displaystyle}
\newcommand{\R}{\mathbb{R}}\newcommand{\N}{\mathbb{N}}\newcommand{\Z}{\mathbb{Z}}
% (un \usepackage{fancyhdr}+en-tête est possible ; il est ignoré côté web)
% =========================================================================
\begin{document}
\begin{corrige}[à quelle distance retrouve-t-on le champ terrestre ?]
On isole $d$ dans $B=\dfrac{2\cdot10^{-7} I}{d}$ :
\[ d=\frac{2\cdot10^{-7}\,I}{B_T}=\frac{2\cdot10^{-7}\times100}{5\cdot10^{-5}}
    =\frac{2\cdot10^{-5}}{5\cdot10^{-5}}=0{,}4\ \si{\metre}=\qty{40}{\centi\metre}. \]
À \qty{40}{\centi\metre} du fil, son champ vaut déjà autant que celui de la Terre.
\end{corrige}

\begin{corrige}[quel courant pour un champ donné ?]
On isole $I$ dans $B=\dfrac{2\cdot10^{-7} I}{d}$, avec $d=\qty{0.02}{\metre}$ :
\[ I=\frac{B\,d}{2\cdot10^{-7}}=\frac{1\cdot10^{-4}\times0{,}02}{2\cdot10^{-7}}
    =\frac{2\cdot10^{-6}}{2\cdot10^{-7}}=\qty{10}{\ampere}. \]
\end{corrige}

\begin{corrige}[bobine plate de $N$ spires]
\begin{enumerate}[label=\alph*)]
  \item Les $N$ spires jointives ajoutent leurs champs : $B=\dfrac{N\,\mu_0 I}{2r}$.
  \item Application, avec $r=\qty{0.05}{\metre}$ :
        \[ B=\frac{50\times4\pi\times10^{-7}\times2}{2\times0{,}05}
            =\frac{50\times1{,}26\cdot10^{-6}\times2}{0{,}1}\approx\qty{1.3e-3}{\tesla}. \]
\end{enumerate}
\end{corrige}

\begin{corrige}[dimensionner un solénoïde]
On isole $N$ dans $B=\dfrac{\mu_0 N I}{L}$ :
\[ N=\frac{B\,L}{\mu_0 I}=\frac{1\cdot10^{-2}\times0{,}5}{1{,}26\cdot10^{-6}\times5}
    =\frac{5\cdot10^{-3}}{6{,}3\cdot10^{-6}}\approx 800\ \text{spires}. \]
\end{corrige}

\begin{corrige}[deux bobines, même densité de spires]
Le champ interne ne dépend que de la densité $n=N/L$ (et de $I$) : $B=\mu_0\,\dfrac{N}{L}\,I=\mu_0 n I$.
\begin{enumerate}[label=\alph*)]
  \item $n_1=\dfrac{100}{0{,}04}=\qty{2500}{\per\metre}$ ; \quad
        $n_2=\dfrac{150}{0{,}06}=\qty{2500}{\per\metre}$.
  \item Les deux densités sont \textbf{égales} et le courant est le même : les deux solénoïdes créent
        donc \textbf{exactement le même champ} interne. Avoir plus de spires ($N_2>N_1$) ne suffit
        pas ; c'est le rapport $N/L$ qui compte.
\end{enumerate}
\end{corrige}

\end{document}
